\chapter{Аудит кандидатов конечного Dirac-seed
\texorpdfstring{$R_2$}{R2} на \texorpdfstring{$H_{15}$}{H15}}
\label{ch:v10-h15-r2-dirac-seed-audit}

Обобщённая флуктуация сохраняет опору исходного конечного оператора.
Поэтому теперь проверяется более ранний вопрос: содержит ли какой-либо
унаследованный оператор два блока
\begin{equation}
 L_L-u_R,qquad Q_L-e_R,qquad
 R_2\sim(\mathbf3,\mathbf2)_{7/6}?
\end{equation}

Введём Hadamard-проектор $\Pi_{R_2}$ на эти два блока и их сопряжения.
Прямое вычисление для пяти внутренних источников даёт
\begin{align}
 \rank\Pi_{R_2}(D_F)
 &=\rank\Pi_{R_2}(A_{\rm Higgs})
 =\rank\Pi_{R_2}(A_{(2)})=0,\nonumber\\
 \rank\Pi_{R_2}(I_{H_{15}})
 &=\rank\Pi_{R_2}(L_{H_{15}})=0.
\end{align}
Таким образом, ни стандартный конечный оператор, ни одноформа Хиггса, ни
Callias--$M_4$ tensor amplifier, ни incidence-лапласиан не содержат
требуемой опоры.

У расширений статус различен. Pati--Salam-поле
$\Sigma\sim(\mathbf2_R,\mathbf2_L,\mathbf{15}_4)$ содержит после ветвления
компоненту типа $R_2$, но требует новой алгебры и нового конечного seed.
Clifford-расширение допускает смешанные weak--colour скаляры, однако само
по себе не выбирает здесь точную двухрёберную пару. Отказ от
$S^0$-реальности даёт исторический лептокварковый канал другого типа, а
строгий зеркальный цикл факторизует перенос через новые вершины и не
создаёт элементарного $D_{R_2}$.

Аудит одиннадцати маршрутов по шести критериям имеет полный ранг $6$ и
нулевой pass-vector. Стандартные операторы получают $5/6$, но проваливают
главный критерий --- саму $R_2$-опору. Нормированный target-loaded
$D_{R_2}$ также получает $5/6$, проваливая наследование. Единственный
структурный маршрут с правильным представлением без прямой записи target ---
Pati--Salam $\Sigma$, но он является новой геометрией.

\begin{theorem}[Нулевая внутренняя проекция на $R_2$]
Проекции стандартного $D_F$, его Higgs-одноформы, допущенного $A_{(2)}$,
Callias--$M_4$ усилителя и лапласиана типового графа на двухблочную
$R_2$-опору равны нулю. Точная ненулевая проекция возникает только у явно
вставленного $D_{R_2}$ или после расширения конечной геометрии.
\end{theorem}

\begin{nogocriterion}[Нормировка не создаёт seed]
Положительный квадратичный родитель может условно нормировать коэффициент
уже заданного $D_{R_2}$, но не превращает нулевой inherited map в
ненулевой. Такой контроль закрывает устойчивость амплитуды, а не
происхождение оператора.
\end{nogocriterion}

\begin{protocol}\begin{enumerate}
\item проверенных внутренних операторов: \textbf{$5$};
\item ненулевых внутренних $R_2$-проекций: \textbf{$0/5$};
\item проверенных маршрутов: \textbf{$11$};
\item прошедших маршрутов: \textbf{$0/11$};
\item exact target-loaded $D_{R_2}$: \textbf{да};
\item inherited map к двум коэффициентам: \textbf{rank $0$};
\item ближайшее структурное расширение: \textbf{Pati--Salam $\Sigma$};
\item physical origin: \textbf{$0/3$};
\item следующий гейт: \textbf{типизированное вложение $R_2$ в $\Sigma$}.
\end{enumerate}\end{protocol}

Машинный сертификат сохранён в
\path{s2t_v10_particle_wrinkle_dislocation_callias_equal_charge_twist_m4_h15_r2_dirac_seed_candidate_audit_gate_results.json}.

\begin{thebibliography}{9}
\bibitem{v10r2seed-paschke}
M.~Paschke, F.~Scheck, A.~Sitarz,
\emph{Can (noncommutative) geometry accommodate leptoquarks?},
Physical Review D 59 (1999), 035003; arXiv:hep-th/9709009.
\bibitem{v10r2seed-kurkov}
M.~A.~Kurkov, F.~Lizzi,
\emph{Clifford Structures in Noncommutative Geometry and the Extended Scalar Sector},
Physical Review D 97 (2018), 085024; arXiv:1801.00260.
\end{thebibliography}